% ================================================================
% ==== FILE: content/k_levels/k7.tex
% ================================================================

% ==============================
%  Ontology of Continua — Core
%  K-level Module: K7 (Social Systems)
%  FULL MODULE — FINAL
% ==============================



\paragraph{Canonical tuple projection note.}
Local K-level displays in this module are projections of the canonical public tuple \(K=(\Omega,\partial\Omega,A,\Theta,P,J,C,k,M)\). When a display omits \(M\), reorders components, or uses level-indexed shorthand, the omitted embedding context is inherited from the surrounding level discussion rather than introduced as a competing definition.

\subsubsection{\texorpdfstring{$K_7$}{K_7} Обзор}
\label{sec:k7-overview}

$K_7$ — уровень \emph{социальных континуумов}: систем, состояния,
оси и потенциалы которых возникают из \emph{взаимодействия нескольких $K_6$
когнитивных континуумов}. Это минимальный уровень, на котором:

\begin{itemize}
    \item появляются общие модели и коллективные знания,
    \item коммуникация становится структурным измерением,
    \item сотрудничество и конфликт приобретают характер динамических сил,
    \item формируется социальная идентичность и распределение ролей,
    \item коллективная память и нормы стабилизируют поведение,
    \item групповые циклы определяют временную организацию.
\end{itemize}

$K_7$ не является суммой агентов: это новый континуум со своими
собственными допустимыми состояниями, порогами, потоками и механизмами
смерти.

% ================================================================
\subsubsection{Пространство состояний \texorpdfstring{$\Omega(K_7)$}{\Omega(K_7)}}

Состояние социального континуума включает:

\[
\Omega(K_7)=
\{
G_{\mathrm{soc}},\;
M_{\mathrm{shared}},\;
C_{\mathrm{comm}},\;
R_{\mathrm{roles}},\;
N_{\mathrm{norms}},\;
H_{\mathrm{hist}},\;
S_{\mathrm{struc}},\;
\mu^{(7)}_t,\;
\Sigma_{\mathrm{identity}}
\}.
\]

Компоненты:
\begin{itemize}
    \item $G_{\mathrm{soc}}$ — социальный граф (агенты как вершины, связи как рёбра),
    \item $M_{\mathrm{shared}}$ — общие модели (коллективные убеждения),
    \item $C_{\mathrm{comm}}$ — структура коммуникации,
    \item $R_{\mathrm{roles}}$ — распределение ролей (статусы, функции),
    \item $N_{\mathrm{norms}}$ — нормы и ограничения,
    \item $H_{\mathrm{hist}}$ — коллективная история и нарративы,
    \item $S_{\mathrm{struc}}$ — структурные позиции и модули,
    \item $\mu^{(7)}_t$ — распределение состояний группы во времени,
    \item $\Sigma_{\mathrm{identity}}$ — структуры социальной идентичности.
\end{itemize}

Допустимость ограничена пространством $M_7$ и порогами групповой когерентности.

% ================================================================
\subsubsection{Граница \texorpdfstring{$\partial\Omega(K_7)$}{\partial\Omega(K_7)}}

Граница включает:
\begin{itemize}
    \item коллапс общих моделей,
    \item фрагментацию социального графа,
    \item разрушение коммуникации,
    \item потерю нормативной когерентности,
    \item нестабильность ролей,
    \item коллапс идентичности,
    \item неспособность стабилизировать коллективную память.
\end{itemize}

Переход через $\partial\Omega(K_7)$ разрушает групповую континуальность.

% ================================================================
\subsubsection{Оси \texorpdfstring{$A(K_7)$}{A(K_7)}}

Из ядра модуля социальных систем:

\[
A(K_7)=
\{
A_{\mathrm{comm}},\;
A_{\mathrm{coop}},\;
A_{\mathrm{conf}},\;
A_{\mathrm{norm}},\;
A_{\mathrm{role}},\;
A_{\mathrm{id}},\;
A_{\mathrm{shared-model}}
\}.
\]

Интерпретация:
\begin{itemize}
    \item $A_{\mathrm{comm}}$ — ось коммуникации (каналы, коды),
    \item $A_{\mathrm{coop}}$ — ось кооперации,
    \item $A_{\mathrm{conf}}$ — ось конфликта/конкуренции,
    \item $A_{\mathrm{norm}}$ — ось нормативности,
    \item $A_{\mathrm{role}}$ — ось роли–позиции,
    \item $A_{\mathrm{id}}$ — социальные идентичности,
    \item $A_{\mathrm{shared-model}}$ — коллективные модели/убеждения.
\end{itemize}

Эти оси нельзя свести к когнитивным различиям $K_6$.

% ================================================================
\subsubsection{Потенциалы \texorpdfstring{$P(K_7)$}{P(K_7)}}

Потенциалы, управляющие социальной динамикой:
\[
P(K_7)=
\{
P_{\mathrm{comm}},\;
P_{\mathrm{coop}},\;
P_{\mathrm{conf}},\;
P_{\mathrm{norm}},\;
P_{\mathrm{id}},\;
P_{\mathrm{role}},\;
P_{\mathrm{shared}}
\}.
\]

Значения:
\begin{itemize}
    \item $P_{\mathrm{comm}}$: эффективность коммуникации,
    \item $P_{\mathrm{coop}}$: преимущество кооперации,
    \item $P_{\mathrm{conf}}$: давление конфликта,
    \item $P_{\mathrm{norm}}$: нормативная когерентность,
    \item $P_{\mathrm{id}}$: потенциал стабилизации идентичности,
    \item $P_{\mathrm{role}}$: давление функциональной роли,
    \item $P_{\mathrm{shared}}$: потенциал когерентности общих моделей.
\end{itemize}

% ================================================================
\subsubsection{Пороги \texorpdfstring{$\Theta(K_7)$}{\Theta(K_7)}}

\[
\Theta(K_7)=
\{
\Theta_{\mathrm{comm}},\;
\Theta_{\mathrm{coh}},\;
\Theta_{\mathrm{norm}},\;
\Theta_{\mathrm{id}},\;
\Theta_{\mathrm{role}},\;
\Theta_{\mathrm{shared}},\;
\Theta_{\mathrm{fragment}}
\}.
\]

Интерпретация:
\begin{itemize}
    \item $\Theta_{\mathrm{comm}}$: минимальная полоса пропускания коммуникации,
    \item $\Theta_{\mathrm{coh}}$: порог когерентности групповой идентичности,
    \item $\Theta_{\mathrm{norm}}$: минимальная нормативная согласованность,
    \item $\Theta_{\mathrm{id}}$: порог устойчивости идентичности,
    \item $\Theta_{\mathrm{role}}$: порог стабильности ролей,
    \item $\Theta_{\mathrm{shared}}$: порог когерентности общих моделей,
    \item $\Theta_{\mathrm{fragment}}$: порог фрагментации социального графа.
\end{itemize}

Переход любого из этих порогов может вызвать коллапс $K_7$.

% ================================================================
\subsubsection{Потоки \texorpdfstring{$J(K_7)$}{J(K_7)}}

\begin{itemize}
    \item $J_{\mathrm{comm}}$ — потоки коммуникации,
    \item $J_{\mathrm{coop}}$ — динамика кооперации,
    \item $J_{\mathrm{conf}}$ — эскалация или деэскалация конфликта,
    \item $J_{\mathrm{norm}}$ — потоки обновления норм,
    \item $J_{\mathrm{id}}$ — переходы идентичности,
    \item $J_{\mathrm{role}}$ — перераспределение ролей,
    \item $J_{\mathrm{shared}}$ — обновление общих моделей,
    \item $J_{\mathrm{PE}\to\mathrm{comm}}$: проекция сигналов ошибки
      предсказания из $K_6$ в коммуникацию,
    \item $J_{\mathrm{comm}\to A_{\mathrm{error}}}$: обратная связь социальной
      коммуникации в когнитивные оси предсказания.
\end{itemize}

Устойчивость требует:
\[
\sigma(J(K_7)) < \sigma_{\max}(K_7).
\]

% ================================================================
\subsubsection{Циклы \texorpdfstring{$C(K_7)$}{C(K_7)}}

\[
C(K_7)=
\{
C_{\mathrm{comm}},\;
C_{\mathrm{coop}},\;
C_{\mathrm{conf}},\;
C_{\mathrm{norm}},\;
C_{\mathrm{id}},\;
C_{\mathrm{shared}}
\}.
\]

Описание:
\begin{itemize}
    \item $C_{\mathrm{comm}}$: коммуникация → интерпретация → ответ → обновление,
    \item $C_{\mathrm{coop}}$: циклы кооперации (контуры доверия и усиления),
    \item $C_{\mathrm{conf}}$: циклы разрешения конфликтов,
    \item $C_{\mathrm{norm}}$: создание норм, их соблюдение, пересмотр,
    \item $C_{\mathrm{id}}$: формирование, переговоры и стабилизация идентичности,
    \item $C_{\mathrm{shared}}$: поддержание общих убеждений и коллективной памяти.
\end{itemize}

Каждый цикл формирует социальную временную структуру.

% ================================================================
\subsubsection{Время \texorpdfstring{$\tau(K_7)$}{\tau(K_7)}}

Социальное время возникает из устойчивых циклов:

\[
\tau(K_7)
= \min_j \tau_{C_j},
\qquad
\Pi(C_j) > \Theta_{\mathrm{time}}.
\]

Время $K_7$ обычно медленнее и более инерционное, чем время $K_6$, из-за
коллективной инерции.

% ================================================================
\subsubsection{Континуальность \texorpdfstring{$k(K_7)$}{k(K_7)}}

По модулю социальных систем:

\[
k_7 =
H(\Omega(K_7))
\cdot
\frac{|C_{\max}^{(7)}|}{|C_{\mathrm{all}}|}
\cdot
\frac{|A_{\mathrm{active}}|}{|A_{\max}|}
\cdot
\left(1-\frac{T_7}{\Theta_7}\right)_+
\cdot
\left(1-\frac{\sigma(J)}{\sigma_{\max}}\right).
\]

Где:
\begin{itemize}
    \item $C_{\max}^{(7)}$ — наибольший когерентный социальный компонент,
    \item $T_7$ — структурное напряжение социального континуума,
    \item $\Theta_7$ — доминирующий порог (когерентности или предела общих моделей),
    \item $\sigma(J)$ — волатильность социальных потоков.
\end{itemize}

Сильная фрагментация $\Rightarrow k_7\to 0$.

% ================================================================
\subsubsection{Структурное напряжение \texorpdfstring{$T(K_7)$}{T(K_7)}}

Источники:
\begin{itemize}
    \item сбои коммуникации,
    \item конфликты норм,
    \item нестабильность ролей,
    \item фрагментация общих моделей,
    \item столкновения идентичностей,
    \item быстрые изменения баланса кооперации и конфликта,
    \item чрезмерная концентрация влияния или структурные узкие места.
\end{itemize}

Коллапс при:
\[
T(K_7) > \Theta_{\mathrm{coh}}.
\]

% ================================================================
\subsubsection{Энергия \texorpdfstring{$E(K_7)$}{E(K_7)}}

\[
E(K_7)
=
E_{\mathrm{comm}}
+ E_{\mathrm{coop}}
+ E_{\mathrm{conf}}
+ E_{\mathrm{norm}}
+ E_{\mathrm{id}}
+ E_{\mathrm{shared}}
+ E_{K_6\to K_7}.
\]

Интерпретация:
\begin{itemize}
    \item стоимость коммуникации,
    \item стоимость поддержания кооперации,
    \item энергия эскалации конфликтов,
    \item затраты на обеспечение норм,
    \item поддержание идентичности,
    \item поддержание общих убеждений,
    \item связанная энергия, унаследованная от $K_6$.
\end{itemize}

% ================================================================
\subsubsection{Операторы на \texorpdfstring{$K_7$}{K_7} (\texorpdfstring{$\Psi$}{\Psi}, \texorpdfstring{$\Phi$}{\Phi}, \texorpdfstring{$\Lambda$}{\Lambda}, \texorpdfstring{$U$}{U}, \texorpdfstring{$\Chi$}{\Chi})}

\begin{itemize}
    \item $\Psi_{7\to8}$: рождение институциональных и цивилизационных континуумов,
    \item $\Phi$: эволюция социальных структур и ролей,
    \item $\Lambda$: объединение подсистем в когерентные группы,
    \item $U$: стабилизация норм и идентичностей,
    \item $\Chi$: разветвление на субкультуры, фракции и специализированные группы.
\end{itemize}

% ================================================================
\subsubsection{Процессы на \texorpdfstring{$K_7$}{K_7}}

Типичные:
\begin{itemize}
    \item коммуникация,
    \item кооперация и формирование коалиций,
    \item эскалация и разрешение конфликтов,
    \item формирование, применение и пересмотр норм,
    \item распределение ролей и формирование авторитета,
    \item создание и трансформация идентичности,
    \item формирование общих нарративов и коллективной памяти,
    \item социальное обучение и диффузия.
\end{itemize}

% ================================================================
\subsubsection{Прогнозы для \texorpdfstring{$K_7$}{K_7}}

\begin{enumerate}
    \item Социальные континуумы коллапсируют, когда когерентность общих моделей
      падает ниже $\Theta_{\mathrm{shared}}$.
    \item Теснота коммуникации предсказывает фрагментацию.
    \item Нестабильность норм предшествует коллапсу идентичности.
    \item Высоко модульные группы устойчивее к шуму и конфликту.
    \item Переход к $K_8$ требует устойчивой нормативной и коммуникационной
      стабильности.
\end{enumerate}

% ================================================================
\subsubsection{Эксперименты для \texorpdfstring{$K_7$}{K_7}}

Возможные эмпирические аналоги:
\begin{itemize}
    \item имитации многоагентного обучения с подкреплением,
    \item исследования фрагментации сетей,
    \item эксперименты по распространению норм,
    \item тесты устойчивости коллективной памяти,
    \item эксперименты с возмущением пропускной способности коммуникации,
    \item измерения когерентности идентичности в социальных группах.
\end{itemize}

% ================================================================
\subsubsection{Коллапс и смерть \texorpdfstring{$K_7$}{K_7}}

Смерть наступает, когда:
\[
\Omega(K_7) = \varnothing.
\]

Механизмы:
\begin{itemize}
    \item фрагментация социального графа,
    \item потеря коммуникативной ёмкости,
    \item распад норм,
    \item коллапс идентичности,
    \item разгонные конкурентные динамики,
    \item неспособность стабилизировать общие модели.
\end{itemize}

% ================================================================
\subsubsection{Фальсифицируемость \texorpdfstring{$K_7$}{K_7}}

Теория предсказывает:
\begin{itemize}
    \item существование порогов когерентности,
    \item измеримые эффекты пропускной способности коммуникации,
    \item чёткую связь между фрагментацией и коллапсом,
    \item устойчивые социальные циклы с характерными периодами,
    \item необходимость структур общих моделей для сохранения группы.
\end{itemize}

Опровержение возможно при наблюдении социальной устойчивости на
крупном масштабе без наличия этих структур.

% ================================================================
\subsubsection{Разветвление / Онтологическое положение \texorpdfstring{$K_7$}{K_7}}

Ветвления:
\begin{itemize}
    \item кооперативные и конкурентные общества,
    \item высоконормативные и слабо-нормативные группы,
    \item идентичностно-центричные и роль-центричные структуры,
    \item централизованные и децентрализованные коммуникационные сети.
\end{itemize}

Онтологическая позиция:
\[
K_6 \to K_7 \to K_8.
\]

% ================================================================
\subsubsection{Связь с M-пространствами}

$M_7$ задаёт:
\begin{itemize}
    \item социальные графовые топологии,
    \item диапазон допустимых нормативных систем,
    \item сложность коммуникации,
    \item когерентность идентичности,
    \item структурную модульность,
    \item связь с $M_6$ (когнитивный слой).
\end{itemize}

Совместимость с $M_7$ определяет, может ли существовать социальный континуум.
